\scanpage{369}
\mspage{108}
% archive sheet number: 368

Je ne trouve d'ailleurs aucune flèche directe entre $C$ et l'un des objets $C_\infty$, $H_C$ ; tout ce que je vois \deleted{\unsure{c'est}} dans ce sens, c'est
\diagram{\begin{tikzcd}[ampersand replacement=\&,column sep=large]
C_{\mathbb N^\circ}\arrow[r,"{\substack{\text{foncteur}\\\text{de locali-}\\\text{sation}}}"]\arrow[d] \& C_\infty\\
C \&
\end{tikzcd}}
où la flèche verticale est la flèche de projection de $C_{\mathbb N^\circ}=C\times_T\mathbb N^\circ$ (\unsure{avec|sur} $\mathbb N^\circ\mathrel{\unsure{\overset{\underline t}{\relbar\joinrel\relbar\mkern-2mu\mapstochar}|\overset{\underline t}{\longrightarrow}}}T$). Ceci n'établit une relation \deleted{\unsure{entre}} topologique directe entre $C$ (vu dans $\mathrm{Cat\,hot}$ ou dans $\mathrm{Hot}$ simplement) et $C_\infty$, que si les deux flèches précédentes sont \deleted{\unsure{des Cat}} tout au moins des \unsure{Hot-équivalences|Hoto-équivalences}. Et le cas-type (problématique) où la question de la relation entre $C$ et $C_\infty$ se pose, est donc celui où $T=\mathbb N^\circ$, i.e. celui d'un système inductif quelconque de \deleted{\unsure{(petites)}} \inserted{petites} catégories, indexé par $\mathbb N$, et interprété comme défini par une cat. fibrée $C=C_{\mathbb N^\circ}$ sur $\mathbb N^\circ$. Dans ce cas particulier, on a du moins une flèche canonique
\[
C\longrightarrow C_\infty
\]
sur laquelle on peut se poser la question si c'est
